% Native hgraph authoring; original AimPaper statements are checked verbatim against AimPaper/main.tex.
\section{The extension spectral sequence}

% AimPaper: 2.1 definition 264 271
\begin{definition}[The extension spectral sequence]
\label{paper:2.1}
\source{AimPaper Definition 2.1}
\level{coarse}
(From \cite{LWX126}, Definition 2.1.)
% BEGIN AIMPAPER STATEMENT
\label{def:ess}
  For a map $f: X\to Y$, we define an $f$-extension spectral sequence (denoted $f$-ESS) as the spectral sequence derived from filtering the above chain complex according to the Adams filtration on $\pi_*$. The $E_0$-page of this spectral sequence is isomorphic to the direct sum of $E_\infty$-pages of the Adams spectral sequences of $X$ and $Y$: 
  $$\ESS{f}_0^{s,t} \cong E_\infty^{s,t}(X) \oplus E_\infty^{s,t}(Y) \Longrightarrow \Ker(\pi_*f)\oplus \Coker(\pi_*f)$$
  where the $d_0$-differential is induced by $f$ on the $E_\infty$-pages.
  
  Differentials in this spectral sequence has the form
  $$d_n^{f}: \ESS{f}_n^{s,t}\to \ESS{f}_n^{s+n, t+n}.$$
% END AIMPAPER STATEMENT


% Correspondence: partial.
% Blueprint note: The extension spectral sequence of the same filtered map is constructed under \texttt{IsBounded} hypotheses on the source and target filtrations. A dedicated construction for an arbitrary classical map $f$, together with the full abutment comparison in the paper, is not yet supplied.
% Blueprint note: Formal counterparts: \item \texttt{KIP126.Core.SpectralSequence.FilteredMapExtension.sequence}: \ref{decl:b59a9249659eb910} \item \texttt{KIP126.Core.SpectralSequence.FilteredMapExtension.differentialRelation\_iff}: \ref{decl:fd30464b637165e0}
\lean{KIP126.Core.SpectralSequence.FilteredMapExtension.sequence,KIP126.Core.SpectralSequence.FilteredMapExtension.differentialRelation_iff}
\uses{decl:b59a9249659eb910,decl:fd30464b637165e0}
\sketch
% Blueprint note: Original statement and context: paper/main.tex.html#L264; paper/2412.10879.pdf.
\end{definition}

% AimPaper: 2.6 corollary 327 330
\begin{corollary}[Vanishing below the Adams filtration]
\label{paper:2.6}
\source{AimPaper Corollary 2.6}
\level{coarse}
(From \cite{LWX126}, Corollary 2.6.)
% BEGIN AIMPAPER STATEMENT
If the Adams filtration of $f$ is $k$, then
    $d^f_i=0$ for $i<k$.
% END AIMPAPER STATEMENT


% Correspondence: no public counterpart identified.
% Blueprint note: No public theorem asserting that $\operatorname{AF}(f)=k$ forces every $d_i^f$ with $i<k$ to vanish has been identified. The result \texttt{action\_filtration} provides related support but does not state this conclusion.
% Blueprint note: Original statement and context: paper/main.tex.html#L327; paper/2412.10879.pdf.
\end{corollary}
